\documentclass[12pt]{amsart}

\usepackage[margin=1in]{geometry}
\usepackage{hyperref}
\hypersetup{
    colorlinks=true,
    linkcolor=black,
    urlcolor=cyan,
    pdftitle={Statement on the Use of Generative AI},
}
\urlstyle{same}

\setlength{\parindent}{0pt}
\setlength{\parskip}{0.8em}

\begin{document}
\thispagestyle{empty}

\begin{center}
  {\large\textsc{Statement on the Use of Generative AI}}\\[5pt]
  \small\textsc{Edward A. Varvak, University of Chicago}
\end{center}

\bigskip

\textbf{Research.} No generative AI was used in my first paper, \emph{The
Monodromy of Compact Lagrangian Fibrations}. For my more recent preprint, 
\emph{The Mumford--Tate group of a very general Lagrangian fiber}, 
I used Gemini 3.1 Pro to proofread the manuscript and to help locate references. 
The arguments and proofs are my own, and no part of the paper was written by AI.

\textbf{Application materials.} I wrote my research statement and publication
list entirely myself, and used AI only to check them for typos. My teaching
statement, cover letter, and generative AI statement were partially written by 
Claude Opus 5.5. I edited all three, and I take full responsibility for all of 
their contents. 

\end{document}
